% --- Template for thesis / report with tktltiki2 class ---
% 
% last updated 2013/02/15 for tkltiki2 v1.02
\documentclass[finnish]{tktltiki2}
%\usepackage[finnish]{babel}
%\usepackage[utf8x]{inputenc}

% tktltiki2 automatically loads babel, so you can simply
% give the language parameter (e.g. finnish, swedish, english, british) as
% a parameter for the class: \documentclass[finnish]{tktltiki2}.
% The information on title and abstract is generated automatically depending on
% the language, see below if you need to change any of these manually.
% 
% Class options:
% - grading                 -- Print labels for grading information on the front page.
% - disablelastpagecounter  -- Disables the automatic generation of page number information
%                              in the abstract. See also \numberofpagesinformation{} command below.
%
% The class also respects the following options of article class:
%   10pt, 11pt, 12pt, final, draft, oneside, twoside,
%   openright, openany, onecolumn, twocolumn, leqno, fleqn
%
% The default font size is 11pt. The paper size used is A4, other sizes are not supported.
%
% rubber: module pdftex

% --- General packages ---
\usepackage[utf8]{inputenc}
\usepackage[T1]{fontenc}
\usepackage{lmodern}
\usepackage{microtype}
\usepackage{amsfonts,amsmath,amssymb,amsthm,booktabs,color,enumitem,graphicx}
\usepackage[pdftex,hidelinks]{hyperref}
\usepackage{apacite}
\usepackage{hyperref}

% Automatically set the PDF metadata fields
\makeatletter
\AtBeginDocument{\hypersetup{pdftitle = {\@title}, pdfauthor = {\@author}}}
\makeatother

% --- Language-related settings ---
%
% these should be modified according to your language

% babelbib for non-english bibliography using bibtex
\usepackage[fixlanguage]{babelbib}
\selectbiblanguage{finnish}

% add bibliography to the table of contents
\usepackage[nottoc]{tocbibind}
% tocbibind renames the bibliography, use the following to change it back
\settocbibname{Lähteet}

% --- Theorem environment definitions ---

\newtheorem{lau}{Lause}
\newtheorem{lem}[lau]{Lemma}
\newtheorem{kor}[lau]{Korollaari}

\theoremstyle{definition}
\newtheorem{maar}[lau]{Määritelmä}
\newtheorem{ong}{Ongelma}
\newtheorem{alg}[lau]{Algoritmi}
\newtheorem{esim}[lau]{Esimerkki}

\theoremstyle{remark}
\newtheorem*{huom}{Huomautus}

% --- tktltiki2 options ---
%
% The following commands define the information used to generate title and
% abstract pages. The following entries should be always specified:

\title{AngularJS web-sovelluskehys}
\author{Eveliina Pakarinen}
\date{\today}
\level{Ohjelmistoarkkitehtuurien harjoitustyö}
\abstract{Tiivistelmä}

% The following can be used to specify keywords and classification of the paper:

\keywords{avain, sanat}

% classification according to ACM Computing Classification System (http://www.acm.org/about/class/)
% This is probably mostly relevant for computer scientists
% uncomment the following; contents of \classification will be printed under the abstract with a title
% "ACM Computing Classification System (CCS):"
% \classification{}

% If the automatic page number counting is not working as desired in your case,
% uncomment the following to manually set the number of pages displayed in the abstract page:
%
% \numberofpagesinformation{16 sivua + 10 sivua liitteissä}
%
% If you are not a computer scientist, you will want to uncomment the following by hand and specify
% your department, faculty and subject by hand:
%
% \faculty{Matemaattis-luonnontieteellinen}
% \department{Tietojenkäsittelytieteen laitos}
% \subject{Tietojenkäsittelytiede}
%
% If you are not from the University of Helsinki, then you will most likely want to set these also:
%
% \university{Helsingin Yliopisto}
% \universitylong{HELSINGIN YLIOPISTO --- HELSINGFORS UNIVERSITET --- UNIVERSITY OF HELSINKI} % displayed on the top of the abstract page
% \city{Helsinki}
%
\usepackage{float}
\usepackage{graphicx}

\addtolength{\oddsidemargin}{-.5in}
\addtolength{\evensidemargin}{-.5in}
\addtolength{\textwidth}{1in}
\addtolength{\topmargin}{-.5in}
\addtolength{\textheight}{1in}
\numberwithin{figure}{section}

\begin{document}

% --- Front matter ---

\frontmatter      % roman page numbering for front matter

\maketitle        % title page
\makeabstract     % abstract page

\tableofcontents  % table of contents

% --- Main matter ---

\mainmatter       % clear page, start arabic page numbering

\setlength{\parindent}{2.5em}
\setlength{\parskip}{0em}
\renewcommand{\baselinestretch}{1.5}
\large

\section{Johdanto}

Johdantoa

% \newpage

\section{Kehyksen esittely}

Esikappale kappale

\subsection{Arkkitehtuurityylit}

Arkkitehtuurityylit 

\subsection{Komponenttien roolijako}

Komponenttien roolijako

\section{Kehyksen arkkitehtuuri}

Kehyksen arkkitehtuuri

\subsection{Yleisarkkitehtuuri}

Yleisarkkitehtuuri

\subsection{Keskeisimmät variaatiopisteet}

Keskeisimmät variaatiopisteet

\section{Kehyksen arviointi}

Kehyksen arviointi

\subsection{Laatuskenaariot}

Laatuskenaariot
\cite{GameTestingAllInOne2005}

\section{Yhteenveto}

Yhteenveto



% Viitataan kuvaan \ref{fig:kaavio5}.
% \begin{figure}[H]
%   \centering
%   \includegraphics[width=14cm]{images/kaavio5}
%   \caption{Kuvateksti}
%   \label{fig:kaavio5}
% \end{figure}


% --- References ---
%
% bibtex is used to generate the bibliography. The babplain style
% will generate numeric references (e.g. [1]) appropriate for theoretical
% computer science. If you need alphanumeric references (e.g [Tur90]), use
%
% \bibliographystyle{babalpha-lf}
%
% instead.


%\bibliographystyle{babplain-lf}
% \newpage
\bibliographystyle{apacite}
\renewcommand{\BRetrieved}[1]{Tarkistettu {#1}, saatavilla\ }
\renewcommand{\BRetrievedFrom}{Tarkistettu saatavilla\ }
\bibliography{references-fi}

% --- Appendices ---

% uncomment the following

% \newpage
% \appendix
% 
% \section{Esimerkkiliite}

\end{document}
